% ================================================
% 文件: 04_半导体器件.tex
% 章节: 半导体器件
% 所属部分: 电子元器件与材料
% 版本: v1.0
% ================================================

\section{半导体器件（二极管、三极管）}
\label{sec:semiconductor_devices}

半导体器件是以半导体材料为基础制成的电子元件。
所谓半导体，是指导电能力介于导体与绝缘体之间的一类固体材料，其电阻率
大致在 $10^{-5}$ 至 $10^{7}\ \Omega\cdot\mathrm{cm}$ 之间。与金属导体
不同，半导体的导电能力对温度、光照、掺杂浓度和外加电场极为敏感，正是
这种可控性使它成为制造放大器件与开关器件的理想材料。最常用的半导体
材料是硅（Si）和锗（Ge），两者均为四价元素；此外还有砷化镓（GaAs）、
氮化镓（GaN）等化合物半导体，它们的电子迁移率更高、耐压更强，广泛
用于微波与大功率器件。

\textbf{本征半导体与掺杂半导体}：完全纯净且晶格完整的半导体称为本征
半导体。在绝对零度时，其价电子全部被共价键束缚，材料呈绝缘状态；
温度升高后部分价电子获得能量跃过禁带成为自由电子，同时在原处留下
带正电的空穴，二者成对出现，称为本征激发。硅的禁带宽度约
$1.12\ \mathrm{eV}$，锗约 $0.67\ \mathrm{eV}$，故锗器件的反向漏电流
和温度敏感性都大于硅器件，这也是硅器件后来完全取代锗器件的主要原因
之一。在本征半导体中掺入微量五价元素（如磷、砷）可形成以自由电子为
多数载流子的 N 型半导体；掺入微量三价元素（如硼、铝）则形成以空穴为
多数载流子的 P 型半导体。掺杂浓度虽极低，却能使载流子浓度提高几个
数量级，从而把电阻率降低几个数量级。在热平衡状态下，电子浓度 $n$ 与
空穴浓度 $p$ 满足质量作用定律
\[
n\,p = n_i^2
\]
其中 $n_i$ 为本征载流子浓度。可见多数载流子浓度提高时，少数载流子
浓度必然相应下降。

\textbf{PN 结的形成与单向导电性}：把 P 型区与 N 型区在同一晶体内相邻
制作，界面处便形成 PN 结。由于两侧载流子浓度悬殊，电子由 N 区向 P 区
扩散、空穴由 P 区向 N 区扩散，在界面附近留下不能移动的杂质离子，
形成耗尽层与内建电场。内建电势差为
\[
V_D = \frac{kT}{q}\ln\frac{N_A N_D}{n_i^2}
\]
其中 $N_A$、$N_D$ 分别为受主与施主掺杂浓度。内建电场的方向恰好阻止
多数载流子继续扩散，最终达到动态平衡。当外加正向电压（P 区接正、
N 区接负）时，外电场削弱内建电场，耗尽层变窄，多数载流子得以大量
越过结面，形成较大的正向电流；当外加反向电压时，耗尽层加宽，只有
少数载流子形成极小的反向饱和电流。这就是 PN 结的单向导电性，也是
一切半导体器件的物理核心。

PN 结的伏安特性可用肖克利方程定量描述：
\begin{equation}
I = I_S\left(e^{\,qU/(nkT)} - 1\right)
\end{equation}
式中 $I_S$ 为反向饱和电流，$q$ 为电子电荷量，$k$ 为玻尔兹曼常数，
$T$ 为绝对温度，$n$ 为理想因子。常温（约 $300\ \mathrm{K}$）下
热电压 $U_T = kT/q \approx 26\ \mathrm{mV}$。由该式可知，正向电压
每增加约 $60\ \mathrm{mV}$，电流增大约十倍，因此正向导通后电压变化
很小：硅二极管的正向导通压降约 $0.6\sim0.7\ \mathrm{V}$，锗二极管
约 $0.2\sim0.3\ \mathrm{V}$，肖特基二极管约 $0.3\sim0.4\ \mathrm{V}$。
反向饱和电流 $I_S$ 随温度上升近似每升高 $10\ ^{\circ}\mathrm{C}$
增大一倍，正向压降则约以 $-2\ \mathrm{mV}/^{\circ}\mathrm{C}$ 的
速率下降，这两条温度规律既是温度补偿电路的依据，也是功率器件热失控
的根源。

\begin{itemize}
    \item \textbf{二极管}：具有单向导电性，由PN结构成。主要用于整流、检波等
    
    \item \textbf{三极管（BJT）}：具有电流放大作用，有NPN和PNP两种类型。共射极放大倍数$\beta = \frac{I_c}{I_b}$
    
    \item \textbf{场效应管（FET）}：电压控制器件，输入阻抗高。分为JFET和MOSFET
\end{itemize}

\textbf{二极管的主要参数与类型}：使用二极管时最需关注四个参数。
最大整流电流 $I_F$ 表示长期正向工作允许的平均电流，超过则因结温过高
而烧毁；最高反向工作电压 $U_{RM}$ 通常取击穿电压的一半以留有余量；
反向电流 $I_R$ 越小说明单向导电性越好；最高工作频率 $f_M$ 受结电容
限制，频率过高时结电容的容抗变小，二极管将失去单向导电作用。按用途
可将二极管分为若干类：整流二极管用于把交流变为脉动直流，结面积大、
结电容大，只能工作于低频；检波二极管用于从调幅波中取出音频包络，
要求结电容小、正向压降低，常用点接触式锗管；开关二极管强调反向恢复
时间短，用于数字电路与脉冲电路；稳压二极管（齐纳二极管）工作在反向
击穿区，利用击穿后端电压几乎不变的特性提供基准电压，其动态电阻
$r_z = \Delta U_Z/\Delta I_Z$ 越小则稳压性能越好，使用时必须串联限流
电阻，使工作电流处于 $I_{Z\min}$ 与 $I_{Z\max}$ 之间；变容二极管利用
反向偏压改变耗尽层宽度从而改变结电容，其容量近似满足
$C_j = C_{j0}/(1 + U_R/V_D)^{m}$，是频率合成器与自动频率控制电路中
实现电调谐的关键器件；肖特基二极管用金属与半导体接触代替 PN 结，
无少子存储效应，反向恢复时间可小于纳秒级，正向压降低，适用于开关
电源整流与高频检波；发光二极管把电能转换为光能，光电二极管则反之，
是光纤通信收发模块的基础。

\begin{table}[htbp]
\centering
\caption{常用二极管类型的特性与典型用途}
\begin{tabular}{|p{2.4cm}|p{3.0cm}|p{3.0cm}|p{3.2cm}|}
\hline
\textbf{类型} & \textbf{工作区域} & \textbf{关键特性} & \textbf{典型用途} \\
\hline
整流二极管 & 正向导通区 & 结面积大，$I_F$ 大，频率低 & 电源整流，把交流变脉动直流 \\
\hline
检波二极管 & 正向小信号区 & 结电容小，正向压降低 & 调幅波检波、包络提取 \\
\hline
开关二极管 & 导通/截止切换 & 反向恢复时间极短 & 数字门电路、脉冲整形 \\
\hline
稳压二极管 & 反向击穿区 & 击穿电压稳定，$r_z$ 小 & 基准电压、过压保护 \\
\hline
变容二极管 & 反向偏置区 & 结电容随反压变化 & 电调谐、频率合成、调频 \\
\hline
肖特基二极管 & 正向导通区 & 压降低，无少子存储 & 开关电源整流、微波混频 \\
\hline
\end{tabular}
\end{table}

\textbf{整流与滤波电路}：二极管最基本的应用是整流。半波整流只利用交流
的半个周期，输出平均电压为 $U_o \approx 0.45U_2$（$U_2$ 为变压器次级
有效值）；桥式全波整流利用两个半周，输出平均电压为
$U_o \approx 0.9U_2$，且脉动频率为电源频率的两倍，更易滤波。整流后
接电容滤波时，输出接近 $\sqrt{2}U_2$，但二极管须承受较大的浪涌充电
电流，其反向峰值电压应按 $\sqrt{2}U_2$（桥式）或 $2\sqrt{2}U_2$
（带中心抽头全波）选取。电容滤波的纹波峰峰值近似为
\[
\Delta U \approx \frac{I_o}{f_r C}
\]
其中 $I_o$ 为负载电流，$f_r$ 为整流后的脉动频率。若纹波仍不满足要求，
可再接 LC 或 RC 滤波网络，或采用稳压电路。

\textbf{双极型晶体管的结构与放大原理}：双极型晶体管（BJT）由两个背靠
背的 PN 结构成，形成发射区、基区、集电区三个区域和发射结、集电结两个
结面，按掺杂顺序分为 NPN 与 PNP 两种类型。其能够放大的结构条件有三：
发射区重掺杂以提供充足的载流子；基区极薄且掺杂很低，使注入的载流子
在其中来不及复合便到达集电结；集电区面积大、掺杂适中，便于收集载流子
并散发功耗。外部条件是发射结正向偏置、集电结反向偏置，即工作于放大区。
此时发射区向基区注入的载流子中，仅有极小部分在基区复合形成基极电流，
绝大部分被集电结的强电场收集形成集电极电流，于是微小的基极电流变化
便控制了很大的集电极电流变化，这就是电流放大的实质。

三个电极电流满足 $I_e = I_b + I_c$。定义共射极直流电流放大系数为
\[
\beta = \frac{I_c}{I_b}
\]
共基极电流放大系数为 $\alpha = I_c/I_e$，两者关系为
\begin{equation}
\alpha = \frac{\beta}{1+\beta}, \qquad \beta = \frac{\alpha}{1-\alpha}
\end{equation}
$\beta$ 一般在数十至数百之间，$\alpha$ 则略小于 1。小信号交流放大
系数记为 $h_{fe}$ 或 $\beta_{ac} = \Delta I_c/\Delta I_b$，数值与
$\beta$ 相近但并不相等。必须注意 $\beta$ 的分散性很大，同型号管子
可相差一倍以上，且随温度升高而增大，因此放大电路必须采用负反馈
（如发射极电阻、分压式偏置）来稳定静态工作点，不能依赖 $\beta$ 的
标称值。

\textbf{三极管的三个工作区}：以集电极电流 $I_c$ 为纵坐标、集电极
与发射极间电压 $U_{ce}$ 为横坐标，以基极电流 $I_b$ 为参变量所得的
输出特性曲线族，可划分为三个区域。放大区中曲线近似水平，
$I_c \approx \beta I_b$，集电极电流几乎只受基极电流控制而与 $U_{ce}$
无关，此区用于线性放大。饱和区位于 $U_{ce}$ 很小处，两个结均正偏，
$I_c$ 不再随 $I_b$ 增加，饱和压降 $U_{ces}$ 约
$0.1\sim0.3\ \mathrm{V}$，相当于开关闭合。截止区中 $I_b \approx 0$，
$I_c$ 仅为极小的穿透电流 $I_{ceo}$，相当于开关断开。数字电路与
功率开关电路正是让晶体管在饱和与截止之间跳变，以获得最低的功率损耗。

晶体管的极限参数决定其安全工作范围：集电极最大允许电流 $I_{cm}$，
超过后 $\beta$ 显著下降；集电极最大允许耗散功率 $P_{cm}$，由
$P_c = U_{ce}I_c$ 计算，必须结合散热条件降额，结温满足
\[
T_j = T_a + P_c\,R_{th}
\]
其中 $R_{th}$ 为总热阻，硅管结温上限通常为
$150\ ^{\circ}\mathrm{C}$，锗管仅约 $85\ ^{\circ}\mathrm{C}$；
以及集电极--发射极击穿电压 $U_{ceo}$。三者围成的区域即安全工作区，
超出任一边界都可能造成永久损坏。此外，特征频率 $f_T$（$\beta$ 下降
到 1 时的频率）与截止频率 $f_\beta$ 反映高频性能，二者关系近似为
$f_T \approx \beta f_\beta$。

\textbf{三种基本组态}：按输入、输出的公共端不同，晶体管放大电路有
三种组态。共射极组态既有电流放大又有电压放大，输出与输入反相，
输入阻抗中等、输出阻抗较高，是最常用的电压放大级。共基极组态电流
放大系数略小于 1 但电压增益高，输入阻抗低、高频特性好、无相位反转，
常用于射频放大与宽带放大。共集电极组态（射极跟随器）电压增益接近 1、
电流增益大，输入阻抗高而输出阻抗低，用于阻抗变换与缓冲隔离。三者的
选用取决于对增益、阻抗与带宽的具体要求。

\textbf{场效应管的工作机理}：场效应管（FET）是电压控制器件，靠栅极
电压产生的电场改变导电沟道的宽窄，从而控制漏源电流。由于栅极几乎
不取电流，其输入阻抗极高：结型场效应管（JFET）可达
$10^{8}\ \Omega$ 以上，绝缘栅型（MOSFET）可达 $10^{12}\ \Omega$
以上，这使它非常适合作为高阻抗输入级、模拟开关和低噪声前置放大器。
FET 只依靠一种载流子导电，故又称单极型晶体管，其温度稳定性优于
BJT，且不存在二次击穿，多管并联时因跨导的负温度系数而能自动均流。

JFET 依靠反向偏置的栅结耗尽层挤压沟道，只能工作在耗尽方式。其在
饱和区的转移特性近似为平方律关系
\begin{equation}
I_D = I_{DSS}\left(1 - \frac{U_{GS}}{U_P}\right)^{2}
\end{equation}
其中 $I_{DSS}$ 为栅源短路时的饱和漏电流，$U_P$ 为夹断电压。
MOSFET 则以二氧化硅薄层把栅极与沟道绝缘，按沟道载流子类型分为
N 沟道与 P 沟道，按零栅压下沟道是否存在分为增强型（须加栅压才形成
沟道）与耗尽型（零栅压即有沟道）。增强型 MOSFET 在饱和区的电流为
\[
I_D = \frac{1}{2}K\left(U_{GS} - U_{T}\right)^{2}
\]
$U_T$ 为开启电压。衡量 FET 放大能力的参数是跨导
$g_m = \Delta I_D/\Delta U_{GS}$，单位为西门子（S）或
$\mathrm{mA/V}$，它对应于 BJT 中的 $\beta$，但量纲不同：前者把
电压变化转换为电流变化，后者把电流变化转换为电流变化。功率
MOSFET 还有一个重要参数是导通电阻 $R_{DS(on)}$，导通损耗为
$P = I_D^2 R_{DS(on)}$，其值随温度升高而增大。

由于栅极氧化层极薄，MOSFET 的栅源击穿电压通常仅
$\pm20\ \mathrm{V}$ 左右，静电放电极易造成击穿。因此存放时必须
将三个电极短接，焊接时应使用防静电烙铁并佩戴接地手环，电路中
则常在栅源之间反向并联稳压二极管加以保护。

\begin{table}[htbp]
\centering
\caption{双极型晶体管与场效应管的对比}
\begin{tabular}{|p{2.6cm}|p{4.0cm}|p{4.4cm}|}
\hline
\textbf{比较项目} & \textbf{双极型晶体管（BJT）} & \textbf{场效应管（FET）} \\
\hline
控制方式 & 电流控制，$I_c = \beta I_b$ & 电压控制，$I_d = g_m U_{gs}$ \\
\hline
参与导电的载流子 & 电子与空穴（双极型） & 仅多数载流子（单极型） \\
\hline
输入阻抗 & 较低，一般数百 $\Omega$ 至数十 k$\Omega$ & 极高，JFET 约 $10^{8}\ \Omega$，MOSFET 可达 $10^{12}\ \Omega$ \\
\hline
主要放大参数 & $\beta$（$h_{fe}$），分散性大 & 跨导 $g_m$，单位 mA/V \\
\hline
温度特性 & $\beta$ 随温度上升，易热失控 & 沟道电阻正温度系数，易于并联均流 \\
\hline
噪声与静电 & 噪声较大，抗静电能力强 & 噪声低，栅极极易被静电击穿 \\
\hline
典型用途 & 音频与射频放大、开关驱动 & 高阻输入级、模拟开关、开关电源 \\
\hline
\end{tabular}
\end{table}

\textbf{型号识别与检测}：国产半导体器件型号一般由五部分组成，
依次表示电极数目、材料与极性、类别、序号和规格。例如 3AD 表示
锗 PNP 低频大功率管，3DG 表示硅 NPN 高频小功率管，2CW 表示硅
稳压二极管，2CP 表示硅整流二极管。使用万用表检测时，可利用
PN 结的单向导电性判别管脚与好坏：二极管正向电阻小、反向电阻大，
两个方向均小说明已短路，均大说明已开路；三极管则可视为两个
背靠背的二极管，通过测量各极间正反向电阻确定基极与管型，
再用 $h_{FE}$ 挡估测放大倍数。对于 MOSFET，因栅极绝缘，
栅源之间正反向电阻均应极大，测量前须先对栅极放电以消除
残余电荷造成的误判。

\textbf{在通信设备中的地位}：半导体器件以体积小、耗电省、
寿命长、无需预热等优点全面取代了电子管，成为现代通信设备的
基础。二极管承担整流、检波、稳压、混频与调谐；晶体管与场效应管
承担从微伏级微弱信号放大到千瓦级射频功率输出的全部有源功能；
而把大量晶体管与无源元件集成在同一芯片上，便形成了下一节讨论的
集成电路。掌握 PN 结的物理特性、器件的极限参数与温度规律，是
正确设计、检测和维修各类通信电路的必要前提。
